% !TEX program = xelatex
% !TEX encoding = UTF-8
% !TEX program = xelatex
% !TEX encoding = UTF-8
\documentclass[xelatex,a4paper,base=11pt,ja=standard,jafont=noto]{bxjsreport}
\setpagelayout{top=2cm,bottom=2cm,left=2cm,right=2cm}

\usepackage{iftex}
\ifXeTeX
  \usepackage{xeCJK}
  \setCJKmainfont{Noto Serif CJK SC}
  \setCJKsansfont{Noto Sans CJK JP}
  \setCJKmonofont{Noto Sans CJK JP}
\else\ifLuaTeX
  \usepackage{luatexja}
  \usepackage{luatexja-fontspec}
  \setmainjfont{Noto Serif CJK SC}
  \setsansjfont{Noto Sans CJK JP}
  \setmonojfont{Noto Sans CJK JP}
\else
  \PackageError{prism}{This document requires XeLaTeX or LuaLaTeX to render Chinese text}{}
\fi\fi

\usepackage{graphicx}
\usepackage[all]{xy}
\usepackage{amsmath, amssymb}
\usepackage{fancybox}
\usepackage{cases}
\usepackage{ascmac}
\usepackage{bm}
\usepackage{theorem}
\usepackage{tikz}
\usepackage{appendix}
\usepackage{lscape}
\usepackage{mdframed} % 枠線を作成するためのパッケージ
\usepackage{titlesec}
\usepackage{booktabs}
\usepackage{setspace}
\usepackage{tikz}
\usepackage{array}
\usepackage{float}
\usepackage{placeins}
\usepackage{geometry}
\geometry{margin=1in}

\titlespacing*{\section}{0pt}{1\baselineskip}{1\baselineskip}
\titlespacing*{\subsection}{0pt}{1\baselineskip}{1\baselineskip}

\renewcommand{\bibname}{参考文献}
\newtheorem{theorem}{Theorem}[section]
\newtheorem{definition}{Definition}[section]
\newtheorem{problem}{Problem}[section]
\newtheorem{example}{Example}[section]
\newtheorem{proof}{Proof}[section]
% Node styles
% Two node styles for game trees: solid and hollow
\tikzset{solid node/.style={circle,draw,inner sep=1.5,fill=black},hollow node/.style={circle,draw,inner sep=1.5}}

\title{経済学入試面接対策ガイドブック}
\author{唯新教育}
\date{2026年1月}

\begin{document}

\maketitle
\tableofcontents

\chapter{背景}
このセクションでは、受験生の基本情報を確認します。回答時間は20～30秒程度で簡潔にまとめましょう。

\section{自己紹介}
初めまして、XXXと申します。出身は中国XXです。XX年に日本に来て、現在XX学校でXXを勉強しています。大学院ではXX（分野）に関する研究をしたいと考えています。本日はよろしくお願いいたします。

\section{学習背景・塾の経験}
大学と塾では経済学の基礎から応用まで学び、特にXX経済学と計量経済学に力を入れました。さらに、自力で教科書や論文を読み、知識を深めました。

\section{留学費用の負担}
主に家族の仕送りで賄っていますが、一部は奨学金とアルバイト収入で負担しています。

\section{言語能力（日本語 vs 英語）}
日本語の方が得意ですが、英語でも論文の読解やディスカッションを行うことは可能です。

\chapter{志望動機（労働経済学の例）}
志望动机/理由除了常见的“为什么来日本”“为什么选择这个学校/大学院”，对难关校来说更多会考察对研究主题的熟练程度，比如自己的研究计划书的概念、模型、参考文献，所选教授的研究内容等。
\section{なぜ日本に来たのか／来日のきっかけは何か}
日本の経済発展と労働市場の特徴に強い関心を持ったことがきっかけです。
特に少子高齢化や非正規雇用の増加といった日本特有の課題を、実証的なデータ分析を通して学びたいと思いました。
また、日本の大学は研究環境が整っており、丁寧な指導と質の高い教育を受けられると考え、留学を決意しました。

\section{なぜこの大学（大学院）を選んだのか}
貴学は実証研究に力を入れており、労働経済学や計量経済学の分野で優れた研究実績を持つ先生方が多い点に魅力を感じました。
また、理論とデータ分析をバランスよく学べるカリキュラムが整っており、自分の研究テーマに最も適した環境だと考えています。
\section{なぜ経済学を学ぼうと思ったのか}
私は高校時代から経済学に強い興味を持っていました。
経済学の知識を活用して経済現象を解釈することで、政策作成や経済発展など、社会に役立つと確信しています。そのため、大学ではXXを専攻し、より深い経済学の知識と分析手法を習得しました。
\par 注：介绍大学专业的时候，如果是经济学相关的直接说即可，如果不是经济学相关的可以介绍自己大学期间学习了经济分析的知识，例如：微分积分，线性代数，概率论，统计学，Stata软件的使用等。
\par 大学卒業後、XX（研究主题）などの問題に気づき、その問題に対応するための方策を探すためにXX経済学（分野名）を研究したいと考えるようになりました。
\section{なぜこの専攻分野を選んだのか}
私は、人々の働き方や所得格差といった社会的課題に関心を持っています。特に、日本では少子高齢化の進行や非正規雇用の増加により、労働市場の構造が大きく変化しています。
こうした現象を理論とデータの両面から分析できるのが労働経済学であり、経済学の中でも最も現実社会に近い分野だと感じました。
将来的には、計量的な手法を用いて労働政策や雇用制度の改善に貢献できる研究を行いたいと考え、この専攻を選びました。
\section{なぜこの教授を選んだのか}
指導を希望している先生は、労働経済学や計量経済学を中心に、所得格差についての実証研究を多く行われています。
先生の研究テーマは、私が関心を持つ「労働市場における格差」や「性別・雇用形態による賃金の違い」と非常に近く、研究アプローチも実証的で明確です。
先生のもとで、理論を踏まえた上で現実のデータを分析し、社会的意義のある成果を出せるよう努力したいと考えています。

\chapter{学習経験}
\section{数学・統計学・計量経済学の習得状況}
\begin{itemize}
    \item \textbf{数学:} 線形代数と微分積分を学び、最適化や確率論も復習しています。
    \item \textbf{統計学:} 平均・分散・回帰分析の基本を学び、PythonやRで練習しています。
    \item \textbf{計量経済学:} OLS、内生性、IV推定、DID分析などを重点的に学習しました。
\end{itemize}
\section{成績表に経済学関連科目がほとんどないが、なぜ経済学を研究したいのか}
大学では他分野を学びましたが、（研究課題）に関心を持つようになり、経済学の理論とデータ分析を通じて原因を明らかにできる点に魅力を感じました。
\section{大学で何を学んだか}
大学では主に社会学と統計学を学び、人々の行動や社会構造の関係を分析しました。その中で数量的な分析方法に興味を持ち、経済学への関心が高まりました。
\section{卒業論文の概要と内容紹介}
回答方法：研究内容+研究方法+数据+结论，可以的话再和研究计划书的内容结合一下
\par 例：
\par 卒業論文では、若者の就業意識と地域経済の関係をテーマにしました。アンケートデータを分析し、都市部と地方での就業意識の違いを統計的に比較しました。
\section{日本語のレベルはどのくらいか／どれくらい学習したか／どのように学んだか}
日本語は大学入学前から学び始め、約3年間になります。日本語学校と独学の両方で勉強し、今年JLPT N１にも合格しました。現在は日常会話や学術的な日本語も問題なく使えるレベルです。
\section{英語のレベルはどのくらいか／どれくらい学習したか／どのように学んだか}
英語は中学から学び、10年以上の学習経験があります、TOEIC ８５０点をとりました。大学では英語文献の読解やレポート作成も行っており、英語論文を読むことにも慣れています。
\section{経済学はどのくらい勉強したか／どのように学んだか}
大学では基礎的なミクロ・マクロ経済を学びました。その後、日本語の教科書やオンライン講座を利用して、理論と計量分析の両方を独学で学び続けています。
\section{数学はどのくらい勉強したか／どのように学んだか}
大学1年から線形代数と微分積分を学び、経済学の勉強と並行して最適化や確率論も復習しています。特に経済モデルを理解するための数学的思考を重視しています。
\section{統計学はどのくらい勉強したか／どのように学んだか}
大学で統計学の授業を履修し、平均・分散・回帰分析などの基本を学びました。その後、PythonやRを使ってデータ分析の練習を行っています。
\section{計量経済学はどのくらい勉強したか／どのように学んだか}
大学卒業後に独学で学び始め、OLSやガウス＝マルコフ定理、内生性、IV推定、DID分析などを重点的に学びました。実際のデータを使って回帰分析も行いました。
\section{なぜ大学のGPAが低いのか／なぜこの科目の成績が悪いのか}
当時は日本語力と専門知識の両方を伸ばす必要があり、時間の配分がうまくできませんでした。しかしその経験を通して学習方法を改善し、今では効率的に学べるようになりました。
\section{なぜ大学の経済学（関連科目）の成績が悪いのか}
初期段階では理論よりも計算部分に苦手意識がありましたが、現在は数学と計量の基礎を固め、理解が大きく深まりました。今ではその分野に最も興味を持っています。
\section{経済学の中で最も得意な分野、最も興味がある分野はどこか}
最も興味があるのは労働経済学です。特に賃金格差や労働市場の構造変化を実証的に分析することに関心があります。
\section{大学時代で最も力を入れた科目は何か}
最も力を入れた科目は労働経済学（または研究関連科目）です。この科目では、〇〇の理論を学び、現実の〇〇を説明できる点に強く興味を持ちました。授業だけでなく、自分でも関連論文を読み、〇〇問題（研究相関）についてのテーマにも理解を深めました。この学習経験が現在の〇〇という研究テーマの設定につながっています。

\section{大学に英語の短期コースはあったか／大学の経済学授業は英語で行われたか}
はい、英語の短期集中コースがあり、参加しました。経済学の授業も一部は英語で行われており、その経験が英語文献を読む力の向上につながりました。
\chapter{大学院進学後の計画}
\section{大学院での学習計画・今後2年間をどのように過ごすつもりか？}
⼊学後は、⼀年⽬で中級と上級ミクロ経済学とマクロ経済学、計量経済学をしっかり基礎をつけた上で、⾼度な実証分析をする⼿法を学び、⼆年⽬で修⼠論⽂の執筆をしたいと思います。
\section{入学しましたら研究テーマを変える可能性がありますか。}
はい、可能性はあると考えています。現在は「研究テーマ」に関心を持っていますが、入学後に授業や先生方のご指導、先行研究への理解が深まる中で、研究テーマと研究方法をより適切な形に調整したいと考えています。
\chapter{卒業後の計画}
\section{卒業後について何か考えているか}
\section{日本で就職する予定か、帰国する予定か}
\section{どのような仕事をしたいか／どの業界・どのような会社に行きたいか}
\section{博士課程進学／就職後も学術研究を続けるか／博士課程での研究方向や展望}
\section{なぜ帰国しないのか}
帰国の場合：
私は⼤学院に進学するのは、⾼齢社会に進んだ⽇本を対象として研究して、中国への参考になれる成果を挙げるために、⽇本の⼤学院で修⼠学位を⽬指したいと思います。将来のキャリアは、中国で公務員試験を受けて、公的政策関係部⾨で、特に⾼齢者に対する厚⽣政策と労働市場に関する政策のデザインについての仕事をしたいと思います。
\chapter{研究計画書に関する基本的な質問（労働経済学の例）}

\section{研究計画書の概要}
私の研究テーマは、「XX業界における収入格差と少子高齢化に関する実証分析」です。
少子高齢化が進む中で、労働力の減少や雇用構造の変化が進み、特にXX業界では賃金や待遇に大きな格差が生じています。
本研究では、労働経済学の理論を基に、パネルデータを用いて収入格差の要因を明らかにし、労働市場の構造的課題を分析することを目的としています。
\section{なぜこのテーマ／方向性を研究したいのか}
日本の少子高齢化は急速に進んでいて、特にXX業界のように人手依存度が高い産業では労働供給の不足が深刻です。
その一方で、賃金格差や雇用形態の違いが労働力の確保を難しくしています。
私は、こうした社会的課題をデータ分析を通して客観的に把握し、政策的な解決策を示したいと考え、このテーマを選びました。
\section{研究課題・研究問題について具体的に説明}
主な研究課題は、「XX業界における収入格差は、年齢・性別・雇用形態・地域要因などによってどのように形成されるのか」という点です。
さらに、少子高齢化が進むことで労働市場の賃金構造にどのような変化が生じているかを実証的に検証します。
\section{研究は何と何の関係を扱うのか}
本研究では、収入（被説明変数）と、雇用形態・年齢構成・性別・経験年数・地域別人口構成・産業構造といった説明変数との関係を分析します。
特に少子高齢化を示す指標（高齢者比率や出生率など）を導入し、人口変動が収入格差に与える影響を推定します。
\section{関連文献はどのように探したか、先行研究はどの程度把握しているか}
関連文献は、日本労働研究雑誌などから検索しました。
特に、賃金格差や少子高齢化に関する実証研究（例：田中(2020)、吉田(2022)など）を参考に、
分析手法としてOLS・固定効果モデル・IV推定などの応用例を整理しました。
\section{この問題／研究に対する現在の主流な見解と結論を紹介}
先行研究では、収入格差は「性別」「企業規模」「雇用形態」の違いに起因するという見解が一般的です。
また、少子高齢化が進むことで高齢労働者の増加や若年層の減少が賃金構造を変化させているとも指摘されています。
ただし、これらの要因を同時に考慮した分析はまだ十分ではなく、その点に本研究の意義があります。
\section{具体的な問題意識}
少子高齢化は単なる人口問題ではなく、賃金格差や労働生産性、産業の持続性にも深く関係しています。この点をデータで明確に示すことが、将来的な労働政策にとって重要だと考えています。
\section{研究の意義}
特にXX業界のような労働集約型産業では、収入の低さが若年層の参入を妨げる要因にもなっています。本研究は、労働経済学の理論を応用して、少子高齢化というマクロ的要因と収入格差というミクロ的現象を統合的に分析する点に特徴があります。
その結果は、賃金制度の改善や人材確保政策の立案に直接的な示唆を与えると期待されます
\section{自分の研究と既存研究との違い（数据的日期的更新或模型的创新等）}
既存研究は単年度データによる断面的な分析が多いのに対し、
本研究ではパネルデータを使用して、同一個人・企業の時系列変化を追跡します。
これにより、個人特性を統制した上で格差の変化を正確に測定できる点が新しい視点です。
\section{現状の研究状況や今後の展望}
現在は、総務省や厚生労働省のデータをもとに変数の整理と基礎統計を進めています。
今後は、固定効果モデルやIV推定を用いて、内生性を考慮した分析を行い、
収入格差のメカニズムをより明確にする予定です。
\section{仮説は何か、なぜそのように設定したのか}
主な仮説は、「少子高齢化が進む地域ほど、収入格差が拡大する」というものです。
その理由は、高齢化による労働供給の変化が、業界内の賃金構造に歪みを生じさせると考えられるためです。
また、若年層の減少によって労働交渉力が弱まり、低所得層の増加につながる可能性もあります。
\section{テーマを変えるとしたら、どのような研究をしたいか}
もしテーマを変えるとすれば、高齢化が企業の生産性や労働参加率に与える影響を研究したいです。
これは現在のテーマと密接に関連しており、同じく労働市場の持続性に焦点を当てる研究です。
\section{卒業論文の紹介}
卒業論文では、「若者の就業意識と地域経済の関係」をテーマに研究しました。
アンケートデータを用いて、地域ごとの雇用機会や賃金水準が就業意識にどう影響するかを分析しました。
この経験から、経済的格差と人口構造の関係に強い関心を持つようになりました。
\section{卒業論文と研究計画の関係}
卒業論文では「意識と経済環境の関係」を扱いましたが、大学院ではそれを一歩進めて、「人口動態と賃金格差の関係」をデータ分析で検証する形に発展させました。
つまり、定性的な関心を定量的な実証研究へと深化させたのが現在の研究計画です。
\section{研究対象は異なる地域での現状がどうですか？}
\subsubsection*{先行研究中涉及到的情况：}
直接把先行研究的内容再阐述一次，例如：先行研究では、主にXX国（地域）を対象に分析が行われています。これらの研究では、「研究背景」の進行に伴い、XX業界において「研究対象」が拡大していること、また雇用形態や制度的違いが「研究対象」に影響していることが指摘されています。その结果、「研究背景」が进む地域ほど劳动供给が不稳定になり、「研究対象」が生じやすいという共通した倾向が确认されています。

\subsubsection*{先行研究中没有涉及到的情况：}
说自己熟悉的内容，例如： 例えば中国では、「研究背景」と「研究対象」が同時に进んでおり、XX業界における「研究対象」の構造は日本や欧米とは异なる特徴を持っています。私はこうした中国の制度的背景や現状にも一定の知识があるため、将来的には国際比较的视点から分析を拡张することも可能だと考えています。
\section{データの収集方法はなんですか？}
本研究では、既存の公的統計データおよびパネルデータを主に使用します。
具体的には、「資料１」に加え、「資料２」などの個票データを利用します。
これらのデータは信頼性が高く、「変数」を同時に把握できるため、「研究対象」を実証的に分析するのに適していると考えています。必要に応じて、複数の年度データを統合し、パネルデータとして分析を行います。
\section{先行研究から学んだことは何ですか？}
先行研究からは、モデル選択、データの使い方、そしてテーマ設定の三点を学びました。特に、固定効果モデルやIV推定による内生性への対応、公的パネルデータの活用方法、マクロな人口動態をミクロの収入格差と結びつける視点が、自分の研究計画の基礎になっています。

\chapter{研究計画書に関する応用的な質問}
このセクションは、使用するモデルによって質問の内容も異なるため、以下のモデルごとに分けて、それぞれに特化した質問を設定するのが望ましいです：RDD（回帰不連続デザイン）、DID（差の差分析）、IV（操作変数法）、重回帰分析、PSM（傾向スコアマッチング）。
\section{なぜこのモデルを選んだのか}
本研究では、収入格差の要因を分析するために固定効果モデルを採用しました。
個人や企業に特有の観測できない要因をコントロールできる点で、単純なOLS推定よりも現実的な推定結果が得られると考えています。
また、時間の経過による人口構造の変化を扱う上でも、パネルデータモデルが最適だと判断しました。
\section{他に候補となるモデルはあるか}
候補としてはIV推定（Instrumental Variables Estimation）も検討しています。
特に、説明変数と誤差項の相関が疑われる場合には、IVや2段階最小二乗法（2SLS）を用いて内生性を補正することを考えています。
\section{被説明変数・説明変数はどれか}
被説明変数は労働者の収入（賃金または時給）です。
説明変数には、年齢、性別、雇用形態、経験年数、学歴、地域の高齢化率、出生率、企業規模などを設定しています。
また、地域経済の構造を反映するため、産業別労働需要や高齢者比率も加えています。
\section{データの出所とデータの種類について説明}
主に総務省の「労働力調査」、厚生労働省の「賃金構造基本統計調査」、および日本家計パネル調査（JHPS/KHPS）を使用します。これらは個票データ（パネルデータ）として構成されており、個人・地域・時系列を一体的に分析できる点が強みです。
\section{多重共線性はあるか、ある場合はどのように考慮・処理するか}
高齢化率や地域経済変数などの間で共線性が生じる可能性があります。
そのため、VIF（Variance Inflation Factor）を用いて共線性を検出し、値が高い場合には変数選択を調整します。また、主成分分析などを補助的に使って情報を圧縮することも検討しています。
\section{内生性はあるか、ある場合はどのように考慮・処理するか}
这个问题可以作为所有模型的前提问题去问。
内生性は多くの実証分析において潜在的に存在すると考えています。そのため、まず理論や制度的背景を踏まえ、説明変数と誤差項の相関が生じる可能性を検討します。内生性の主な要因としては、逆の因果関係、観測できない要因の欠落、測定誤差が挙げられます。内生性が疑われる場合、パネルデータが利用できれば固定効果モデルを用いて不変の個人特性を統制します。また、逆の因果関係や同時性が問題となる場合には、操作変数法を用い、2段階最小二乗法で推定します。さらに、OLSや固定効果、IV推定の結果を比較し、推定結果の頑健性を確認することが重要だと考えています。
\section{データ量が不十分または不完全な場合、どのように対応するか}
データの欠損がある場合には、まず欠損パターンを分析し、欠損がランダムであれば平均補完や回帰補完を行います。
一方で、サンプルが小さい場合は、複数年度データを統合して観測数を増やす対応をします。
\section{類似した研究経験はあるか}
卒業論文で、若年層の就業意識と地域経済の関係を分析しました。
その際に基本的な回帰分析を行い、データ整理から推定までの流れを経験しました。
この研究経験が、今回の実証分析の基礎になっています。
\section{計量分析ソフトを使用した経験はあるか／データ分析・回帰分析の経験はあるか}
はい、あります。主にStataとRを使用して回帰分析、相関分析、パネル推定などを行いました。また、Pythonを使ってデータの前処理や可視化も行っており、統計的処理には慣れています。
\section{実証が理論を裏付けることは可能か／理論が実証に置き換えられることは可能か}
はい、可能です。理論モデルは労働市場における均衡を前提にしていますが、実証分析を通して「現実がどこまで理論通りか」を検証できます。また、データ分析の結果から理論モデルを修正・発展させることも重要だと考えています。

\section{モデル内の説明変数は被説明変数を完全に説明できるか}
完全には説明できません。賃金には観測できない要因（努力、職場環境、モチベーションなど）も影響します。そのため、固定効果を導入して個人の不変要因を統制し、モデルの精度を高めています。
\section{クロスセクションデータ・時系列データ・パネルデータの特徴と利点・欠点を説明}
クロスセクションデータは一時点の比較が可能ですが、時間的変化を捉えにくい点が欠点です。時系列データは変化の傾向を分析できますが、個人差を考慮できません。
パネルデータは両者の長所を兼ね備えており、個人や地域の固定効果を扱える点で最も柔軟です。本研究ではこの特徴を活かしています。
\section{誤差項をどのように理解しているか}
誤差項は、モデルで説明できない要因をまとめた部分だと理解しています。
観測できない個人特性や測定誤差、外的ショックなどが含まれます。
この誤差が説明変数と独立であることが、OLS推定における重要な仮定です。
\section{各変数／特定の変数を選んだ理由を説明}
収入を被説明変数に選んだのは、経済的格差を定量的に捉えやすいためです。
雇用形態・年齢・性別・高齢化率などの変数は、収入格差に直接的な影響を与える主要な要因として設定しました。
また、地域の高齢化率や出生率を加えることで、人口動態の変化を経済格差と結びつけて分析することができます。

\chapter{経済学の基礎概念キーワード}
\section*{はじめに}

教授如果想知道你的经济学知识储备量的话，会考察一些比较常见的定义和性质，如果有出题目让你做的话，难度也不会太大（一般刷小黄书和过去问就可以对应），如果教授问到了自己不太熟悉的知识点，也不用太紧张，首先冷静思考一下，然后跟教授说这一个点自己不是很熟，但是自己知道关联的知识点，再补充一下自己打算通过什么样的学习计划来补足知识储备。

\subsection*{注意点}

\begin{enumerate}
\item 即使有不知道的也一定要保持自信，让教授知道你有努力做准备只是刚好这个点不太熟悉。
\item 微观宏观的知识点list可以参考『公務員試験　新スーパー過去問ゼミ　ミクロ経済学』『公務員試験　新スーパー過去問ゼミ　マクロ経済学』的附录。
\end{enumerate}
\newpage
\section{ミクロ経済学}
\begin{table}[H]
\centering
\small
\begin{tabular}{|>{\raggedright\arraybackslash}p{3.5cm}|>{\raggedright\arraybackslash}p{3.5cm}|>{\raggedright\arraybackslash}p{3.5cm}|>{\raggedright\arraybackslash}p{3.5cm}|}
\hline
アドヴァース・セレクション & 異時点間の最適消費 & 一般均衡分析 & 売上高最大化 \\
\hline
エッジワースのボックス・ダイヤグラム & エンゲル曲線 & 応益原則 & オファー・カーブ \\
\hline
外部効果 & 価格硬直性 & 価格消費曲線 & 下級財 \\
\hline
確実性等値 & 加重限界効用均等法則 & 加重限界生産性均等法則 & 課税 \\
\hline
寡占 & 価値財 & 可変費用 & カルテル \\
\hline
カルテル・モデル & 関税 & 間接効用関数 & 完全競争市場 \\
\hline
機会損失（機会費用） & 危険愛好的 & 危険回避的 & 危険中立的 \\
\hline
技術的限界代替率 & 期待効用最大化 & ギッフェン財 & 規模に関して収穫一定 \\
\hline
規模に関して収穫逓減 & 規模に関して収穫増加 & 規模の経済性 & 逆需要曲線 \\
\hline
逆選択 & 供給曲線 & 供給の価格弾力性 & 競争均衡 \\
\hline
共同消費 & 共謀モデル & 均衡の安定性 & クールノー点 \\
\hline
クールノー＝ナッシュ均衡 & クールノー・モデル & 屈折需要曲線 & くもの巣調整メカニズム \\
\hline
クラブ財 & 繰り返しゲーム & 経済的厚生 & 契約曲線 \\
\hline
ゲーム・ツリー & ゲーム理論 & 限界効用 & 限界代替率 \\
\hline
限界収入 & 限界生産性逓減法則 & 限界費用 & 限界費用価格形成原理 \\
\hline
限界変形率 & 限界便益 & 顕示選好 & 交易条件 \\
\hline
公共財 & 厚生経済学の第1定理 & 厚生経済学の第2定理 & 効用 \\
\hline
効用最大化 & コースの定理 & 固定費用 & コブ＝ダグラス型生産関数 \\
\hline
コモンプール財 & 混合戦略 & コンツェルン & 最適消費点 \\
\hline
最適要素投入 & 最適労働供給 & サブゲーム & サブゲーム完全均衡 \\
\hline
差別独占 & サミュエルソン条件 & サンク・コスト & 死荷重 \\
\hline
シグナリング & 支出関数 & 市場の失敗 & 市場メカニズム \\
\hline
自然独占 & 指数法則 & 私的限界費用 & 支配戦略 \\
\hline
支配戦略均衡 & 資本 & 資本の限界生産性 & 資本レンタル価格 \\
\hline
社会的限界費用 & 社会的厚生 & 奢侈財 & 従量税 \\
\hline
囚人のジレンマ & 自由貿易 & シュタッケルベルク・モデル & 需要曲線 \\
\hline
需要の価格弾力性 & 需要の交差弾力性 & 需要の所得弾力性 & 準公共財 \\
\hline
上級財 & 小国の仮定 & 消費者余剰 & 情報の非対称性 \\
\hline
情報の不確実性 & 情報の不完全性 & 所得効果 & 所得消費曲線 \\
\hline
数量調整 & スクリーニング & ストルパー＝サミュエルソン定理 & 静学ゲーム \\
\hline
生産可能性フロンティア & 生産可能領域 & 生産関数 & 生産者余剰 \\
\hline
生産補助金 & 生産要素 & 正常財 & 贅沢財 \\
\hline
\end{tabular}
\end{table}
\begin{table}[H]
\centering
\small
\begin{tabular}{|>{\raggedright\arraybackslash}p{3.5cm}|>{\raggedright\arraybackslash}p{3.5cm}|>{\raggedright\arraybackslash}p{3.5cm}|>{\raggedright\arraybackslash}p{3.5cm}|}
\hline
製品差別化 & 絶対優位 & 線形生産関数 & 操業停止点 \\
\hline
粗代替財 & 粗補完財 & 損益分岐点 & 代替効果 \\
\hline
代替財 & 短期限界費用曲線 & 短期総費用曲線 & 短期平均費用曲線 \\
\hline
超過供給 & 超過供給価格 & 超過需要 & 超過需要価格 \\
\hline
長期均衡点 & 長期限界費用曲線 & 長期総費用曲線 & 長期平均費用曲線 \\
\hline
貯蓄 & 貸金利率 & 点弾力性 & 同一性命題 \\
\hline
動学ゲーム & 等生産量曲線 & 中間値定理 & 道徳的危険 \\
\hline
等比数列 & 等費用線 & 独占 & 独占均衡点 \\
\hline
独占的競争 & 独立財 & トラスト & トリガー戦略 \\
\hline
ナッシュ均衡 & 排出権取引市場 & 排除不可能性 & パレート改善 \\
\hline
パレート効率性 & パレート最適 & 比較生産費説 & 比較優位 \\
\hline
非結合性 & ピグー税 & 必需財 & ヒックス需要曲線 \\
\hline
非排除性 & 微分 & 費用 & 費用最小化 \\
\hline
費用逓減産業 & フォーク定理 & 不確実性 & 不完全競争市場 \\
\hline
複占 & 部分均衡分析 & 部分ゲーム & 部分ゲーム完全均衡 \\
\hline
プライス・テイカー & プライス・メーカー & フリーライダー & フルコスト原理 \\
\hline
プレイヤー & 平均可変費用 & 平均固定費用 & 平均費用 \\
\hline
平均費用価格形成原理 & ヘクシャー＝オリーンの定理 & ベルトラン・モデル & 偏微分 \\
\hline
貿易三角形 & 包絡線 & ボーエン条件 & 補完財 \\
\hline
補償需要曲線 & 補助金 & マークアップ原理 & マーシャル的調整 \\
\hline
マーシャル需要曲線 & 埋没費用 & マクシ・ミン原理 & 無差別曲線 \\
\hline
モニタリング & モラル・ハザード & 輸入関税 & 輸入数量制限 \\
\hline
輸入割当 & 要素価格 & 要素価格均等化定理 & 予算制約線 \\
\hline
余剰 & 余剰損失 & 余剰分析 & ラーナーの独占度 \\
\hline
リカード・モデル & 利潤最大化 & リスク & リスク・プレミアム \\
\hline
利得 & 利得行列表 & リプチンスキーの定理 & 流動性制約 \\
\hline
留保価格 & リンダール均衡 & リンダール・メカニズム & レオンチェフ型生産関数 \\
\hline
レオンチェフの逆説 & 劣等財 & レモンの原理 & 連関財 \\
\hline
労働 & 労働の限界生産性 & 割引現在価格 & ワルラス的調整 \\
\hline
ワルラス法則 & & & \\
\hline
\end{tabular}
\end{table}

\newpage

\section{マクロ経済学}
\begin{table}[H]
\centering
\small
\begin{tabular}{|>{\raggedright\arraybackslash}p{3.5cm}|>{\raggedright\arraybackslash}p{3.5cm}|>{\raggedright\arraybackslash}p{3.5cm}|>{\raggedright\arraybackslash}p{3.5cm}|}
\hline
2段階アプローチ & 6つの定型化された事実 & 45度線分析 & AD-AS分析 \\
\hline
AD曲線 & AS曲線 & BP曲線 & CGPI \\
\hline
CPI & DI & GDE & GDI \\
\hline
GDP & GDPギャップ & GDPデフレーター & GNP \\
\hline
NDP & NI & NNI & SNA \\
\hline
IS-LM分析 & IS曲線 & ISバランス式 & LM曲線 \\
\hline
UV曲線 & 2\%ルール & M1 & M2 \\
\hline
M3 & 安全資産 & 安定成長 & イールド・カーブ \\
\hline
一般物価指数 & インカム・ゲイン & インターバンク市場 & インフレーション \\
\hline
インフレ・ギャップ & インフレ供給曲線 & インフレ需要曲線 & 売りオペ \\
\hline
営業余剰 & オークンの法則 & オープン市場操作 & 買いオペ \\
\hline
海外からの純要素所得受取 & 海外需要 & 外国為替相場 & 外需 \\
\hline
開放経済 & 拡張的財政政策 & 可処分所得 & 加速度原理 \\
\hline
株式 & 株式市場 & 貨幣ヴェール & 貨幣錯覚 \\
\hline
貨幣残高方程式 & 貨幣市場 & 貨幣需要の利子弾力性 & 貨幣数量説 \\
\hline
貨幣供給量 & 貨幣の所得速度 & 貨幣の流通速度 & カルドア \\
\hline
為替レート & 間接金融 & 完全雇用国民所得 & 協調的金融政策 \\
\hline
企業物価指数 & 危険資産 & 基礎消費 & 期待インフレ率 \\
\hline
キチンの波 & キャピタル・ゲイン & 均衡国民所得 & 均衡成長 \\
\hline
銀行の銀行 & 均衡予算乗数定理 & 均斉成長 & 金融資産市場 \\
\hline
金融収支 & 金融政策 & 金利操作 & 金利平価説 \\
\hline
相対所得仮説 & クズネッツ & クズネッツ型消費関数 & クラウディング・アウト \\
\hline
景気循環 & 経常収支 & ケインジアン & ケインズ \\
\hline
ケインズ型消費関数 & 減価償却 & 限界消費性向 & 限界税率 \\
\hline
現金通貨 & 現実成長率 & 衒示的消費 & 建設国債原則 \\
\hline
建築循環 & 公開市場操作 & 広義流動性 & 公債 \\
\hline
公債の市中消化原則 & 恒常所得仮説 & 構造的失業 & 公的需要 \\
\hline
購買力平価説 & 効率市場仮説 & 合理的期待形成学派 & 合理的期待形成仮説 \\
\hline
コール市場 & 国際資本移動 & 国際収支表 & 国内需要 \\
\hline
国内所得 & 国内純生産 & 国内総固定資本形成 & 国内総支出 \\
\hline
国内総所得 & 国内総生産 & 国富 & 国民経済計算体系 \\
\hline
国民純所得 & 国民所得 & 国民総所得 & コストプッシュ・インフレ \\
\hline
固定為替相場制 & 固定資本減耗 & 古典派 & 古典派の第一公準 \\
\hline
古典派の第二公準 & 古典派の二分法 & 雇用者報酬 & コンドラチェフの波 \\
\hline
サービス収支 & 債券 & 債券市場 & 在庫管理アプローチ \\
\hline
在庫品増加 & 最終需要 & 財政政策 & 裁定取引 \\
\hline
サプライサイド学派 & 産業連関表 & 産業連関分析 & 産出増大型技術進歩 \\
\hline
三面等価の原則 & 時間的相対所得仮説 & 資産動機 & 支出面から見たGDP \\
\hline
\end{tabular}
\end{table}
\begin{table}[H]
\centering
\small
\begin{tabular}{|>{\raggedright\arraybackslash}p{3.5cm}|>{\raggedright\arraybackslash}p{3.5cm}|>{\raggedright\arraybackslash}p{3.5cm}|>{\raggedright\arraybackslash}p{3.5cm}|}
\hline
市場価格表示 & 自然失業率 & 自然失業率仮説 & 自然成長率 \\
\hline
実質GDP & 自発的失業 & 資本 & 資本移転等収支 \\
\hline
資本減耗 & 資本ストック調整原理 & 資本増大型技術進歩 & 資本蓄積の黄金律 \\
\hline
習慣形成仮説 & 修正版フィリップス曲線 & ジュグラーの波 & 需要不足失業 \\
\hline
純間接税 & 純輸出 & 証券投資 & 小循環 \\
\hline
乗数効果 & 消費者物価指数 & 所得税率 & ジョルゲンソン型投資関数 \\
\hline
新貨幣数量説 & 新古典派型投資関数 & 新古典派経済成長理論 & 信用乗数 \\
\hline
信用創造 & 信用秩序維持政策 & スタグフレーション & ストック \\
\hline
静学的期待形成仮説 & 政策金利 & 政策割当論 & 成長会計 \\
\hline
政府最終消費支出 & 政府支出乗数 & 政府の銀行 & 絶対所得仮説 \\
\hline
潜在国民所得 & 潜在成長率 & 全要素生産性の上昇率 & 総供給曲線 \\
\hline
操作変数 & 総需要管理政策 & 総需要曲線 & 総需要・総供給分析 \\
\hline
相対所得仮説 & 租税 & 租税乗数 & ソロー＝スワン型モデル \\
\hline
ソロー中立型技術進歩 & 第一次所得収支 & 第二次所得収支 & 短期消費関数 \\
\hline
中央銀行 & 中間目標 & 長期金融市場 & 長期消費関数 \\
\hline
長期波動 & 調整係数 & 直接金融 & 直接投資 \\
\hline
貯蓄 & 貯蓄関数 & 貯蓄のパラドックス & 定額税 \\
\hline
ディマンドプル・インフレ & 適応的期待形成仮説 & デフレ・ギャップ & デモンストレーション効果 \\
\hline
デューゼンベリー & 投機的動機 & 投資 & 投資関数 \\
\hline
投資効果曲線 & 投資の限界効率曲線 & 投資の限界効率理論 & 投資の調整費用 \\
\hline
投資の二重性 & 投資の利子弾力性 & 投入係数 & トービン \\
\hline
トービンのq理論 & 取引動機 & 内需 & 内生的経済成長理論 \\
\hline
ナイフ・エッジ原理 & 日本銀行 & ニュー・ケインジアン & パーシェ指数 \\
\hline
ハイパワード・マネー & 発券銀行 & ハロッド中立型技術進歩 & ハロッド＝ドーマー型経済成長理論 \\
\hline
ピグー効果 & 非自発的失業 & ヒックス中立的技術進歩 & ビルトイン・スタビライザー \\
\hline
不安定性原理 & フィッシャーの交換方程式（数量方程式） & フィッシャー方程式 & フィッシャー効果 \\
\hline
フィリップス曲線 & 封鎖経済 & 付加価値 & 不胎化政策 \\
\hline
物価版フィリップス曲線 & フリードマン & ブルーデンス政策 & フロー \\
\hline
フロート制 & 分配面から見たGDP & 平均消費性向 & 閉鎖経済 \\
\hline
ベヴァレッジ曲線 & 変動為替相場制 & ペンローズ曲線 & 貿易収支 \\
\hline
法定準備率操作 & 保証成長率 & ポリシー・ミックス & 本源的預金 \\
\hline
マーシャルのk & マクロ生産関数 & 摩擦的失業 & マネーサプライ \\
\hline
マネーストック & マネタリー・ベース & マネタリスト & マンデル＝フレミング・モデル \\
\hline
民間最終消費支出 & 民間消費 & 民間投資 & 無限等比数列の和の公式 \\
\hline
名目GDP & メニュー・コスト理論 & モディリアーニ & 有効需要の原理 \\
\hline
輸出 & 輸入 & 輸入関数 & 要素費用表示 \\
\hline
預金通貨 & 予備的動機 & ライフサイクル仮説 & ラスパイレス指数 \\
\hline
ラチェット効果 & 利子率 & リスク・プレミアム・レート & 流動資産仮説 \\
\hline
流動性選好説 & 流動性のわな & 労働供給 & 労働供給曲線 \\
\hline
労働需要 & 労働需要曲線 & 労働増大型技術進歩 & 労働の限界生産性 \\
\hline
割引現在価値 & ワルラス法則 & ルーカス批判 & レーガノミクス \\
\hline
レオンチェフ型生産関数 & & & \\
\hline
\end{tabular}
\end{table}

\newpage
\section{計量経済学}
参考山本勲教授的教科书，把整个计量的知识点用语言表述的方式连接，如果问到某一个知识点的话，就把前后两个知识点的内容一起说即可。
\par \textbf{注意点}：除了这些基础知识之外还需要深入自己的研究计划书所涉模型。
\par 経済データを分析するとき、まず最も基本的な方法は最小二乗法OLSを用いてパラメータを推定することです。OLSを使って推定した結果は、BLUE、つまり「Best Linear Unbiased Estimator（最良線形不偏推定量）」の性質を持っています。
\par このBLUEには主に4つの特徴があります。
Linearity（線形性）、Unbiasedness（不偏性）、Efficiency（効率性）、そしてConsistency（一致性）**です。
これらの性質を保証するためには、ガウス＝マルコフの定理が重要です。
この定理によると、誤差項について3つの仮定が必要です。
\par それは、
\par １、誤差項が均一な分散（等分散性）を持つこと、
\par ２、共分散がゼロであること（誤差同士が独立していること）、
\par ３、説明変数と誤差項が独立であることです。
\par しかし、現実のデータではこれらの仮定が満たされないことも多いです。
\par その場合、それぞれに応じた修正手法を取ります。
\par まず、等分散性が満たされない場合、すなわち誤差の分散が不均一なときには、加重最小二乗法（Weighted Least Squares, WLS）や一般化最小二乗法（Generalized Least Squares, GLS）を用いて修正することができます。
\par 次に、誤差項の共分散がゼロでない、つまり誤差に相関がある場合、特にパネルデータを扱うときには、変量効果モデル（Random Effects Model）を用いてGLSで推定するのが一般的です。
\par また、説明変数と誤差項が独立でない場合、つまり内生性バイアスがあるときには、操作変数法（Instrumental Variables, IV）を用います。
\par このとき、IVとして選ぶ変数には2つの条件が必要です。
\par １、説明変数（X）には影響を与えるが、
\par ２、被説明変数（Y）には直接影響を与えないことです。
\par さらに、パネルデータを使用していて、観測できない要因が内生性を生じさせる場合には、固定効果モデル（Fixed Effects Model）を用いることで、その問題を解決できます。
このように、データの特性や仮定の成否に応じて、適切な推定方法を選ぶことが計量経済学の重要なポイントだと考えています。
\newpage
\section{簡単な数学知識の回答例}
\subsection*{関数}
ある変数の値に対して、一つの値を対応させる関係です。
\subsection*{微分}
関数の変化率を表します。経済学では、限界効果や限界利益を示すために使われます。利潤関数を微分しました、利潤最大化条件も求めます。
\subsection*{偏微分}
複数の変数を持つ関数で、他の変数を固定したまま一つの変数だけで微分することです。経済学的には、他の条件を一定としたときの限界効果を表します。
\chapter{その他の問題}
\section{不合格になりましたどうしますか？}
もし今回ご縁がなかった場合でも、これまで準備してきた研究計画や基礎力をさらに磨き、再挑戦したいと考えています。
特に、研究テーマや分析手法について改めて見直し、より完成度の高い形にして次の機会につなげたいと思います。
ただし、現時点では貴学で学ぶことを第一に考え、全力で準備しています。
\section{他の学校の選考状況はどうですか？}
他大学にも出願はしていますが、研究内容や指導体制を総合的に考えると、貴学が第一志望です。特に、実証研究を重視する教育方針と先生方の研究分野が自分の関心と最も合致していると感じています。そのため、可能であればぜひ貴学で研究に取り組みたいと考えています。

\chapter{逆質問}
\section{個人面}
\begin{itemize}
    \item 入学前に、研究や計量分析の基礎力をさらに高めるために、特に意識して準備しておくべき点があれば教えていただけますでしょうか。
    \item 修士課程を修了した学生の中で、特に評価されている共通点や能力があれば教えていただけますでしょうか。
\end{itemize}
\section{教授面}
\begin{itemize}
    \item 教授先生から見ると合格されました学生さんはどんな品質を持っていますか？
\item 入学後、研究テーマを発展させる際に、先生が特に重視される視点や姿勢があればお聞きしたいです。
\end{itemize}
\section{学校面}
\begin{itemize}
    \item 貴学では、修士課程の学生が研究発表や学会参加を行う機会はどの程度ありますか。
\item 貴学では、修了後の進学や就職に向けて、どのようなキャリア支援が用意されていますか。
\end{itemize}
\end{document}
